% ==========================================================================
\section{Algorithms and baselines}
\label{app:baselines}
% ==========================================================================
All methods start from $\theta=0$ and use full-batch gradients, the step size
$\eta=1/\max_kL_k$, $\tau=10$ local steps per round and $T=150$ rounds. They are
evaluated as in Section~\ref{sec:exp_setup}. The baselines implement their published
update rules.
\begin{itemize}[leftmargin=1.2em,itemsep=2pt]
  \item \textbf{FedAvg}: averaging with $p_k\propto n_k$; this is Algorithm~\ref{alg:ptopofl} with $M=1$ and $w_k=p_k$.
  \item \textbf{FedProx}: local objective $F_k(\theta)+\frac{\mu}{2}\norm{\theta-\theta^t}^2$, $\mu=0.1$, step size $1/(L+\mu)$.
  \item \textbf{SCAFFOLD}: option~II control variates, global step size $1$, full participation, weights $p_k$ \citep{karimireddy2020scaffold}.
  \item \textbf{pFedMe} \citep{t2020personalized}: $\lambda=15$.
    \begin{itemize}[leftmargin=1.2em,itemsep=1pt]
      \item Each of the $\tau$ local iterations approximately solves the proximal problem
        $\min_\theta F_k(\theta)+\frac\lambda2\norm{\theta-w_k}^2$ with $10$ warm-started
        gradient steps of size $1/(L+\lambda)$, then sets
        $w_k\gets w_k-\tfrac12(w_k-\tilde\theta_k)$.
      \item The server averages with $\beta=1$.
      \item A client's personalised model is the proximal point of the final global model.
    \end{itemize}
  \item \textbf{Ditto} \citep{li2021ditto}: a FedAvg global model, plus personal models
    updated with $\tau$ steps per round on $F_k(v)+\frac{\lambda}{2}\norm{v-w^t}^2$,
    $\lambda=0.1$.
  \item \textbf{IFCA} \citep{ghosh2020efficient}: $M=2$.
    \begin{itemize}[leftmargin=1.2em,itemsep=1pt]
      \item The cluster models are initialised from two distinct, randomly chosen clients,
        each trained locally for $20$ steps; identical initialisations would make every
        client choose the same cluster.
      \item In each round, every client selects the cluster model with the lowest training
        loss; each cluster then averages its members' local models.
    \end{itemize}
  \item \textbf{PACFL-style}: each client sends the top-$3$ left singular vectors of its
    (features $\times$ samples) training matrix \citep{vahidian2023efficient}. The
    dissimilarity is the sum of the principal angles, and clustering uses average linkage
    with $M=2$. Training is identical to \ptopofl, with $w_k\propto n_k$ and the same
    calibrated $\alpha$, so the comparison isolates the descriptor.
  \item \textbf{Robust aggregators}: unweighted coordinate-wise median and trimmed mean
    (trimming fraction $0.2$) of the clients' local models \citep{yin2018byzantine}.
  \item \textbf{References.} The random partition is a uniformly random permutation of the
    \ptopofl\ labels, so cluster sizes are identical. The oracle partition is $k$-means
    ($20$ restarts) on the exact local optima, computed by Newton's method. Both use the
    \ptopofl\ training with $w_k\propto n_k$.
\end{itemize}

% ==========================================================================
\section{Descriptor details, implementation check and cost}
\label{app:cost}
% ==========================================================================
\paragraph{Layout.} Coordinates $0$--$19$ hold $\mathrm{BC}_0$, $20$--$39$ hold
$\mathrm{BC}_1$, and coordinates $40$--$45$ hold $(N_0,N_1,E_0,E_1,A_0,A_1)$. Records with
zero pairwise distance, which occur with discrete attributes, produce pairs of zero
persistence; these are excluded from $N_j$, $E_j$, $A_j$ and the Betti curves. Constant
coordinates are left unchanged by the min--max normalisation.

\paragraph{Correctness of $H_1$.} The released code computes $H_1$ with Ripser when it is
installed, and otherwise with a direct $\mathbb{Z}_2$ boundary-matrix reduction of the same
truncated triangle filtration. On random point clouds the two give identical pairs, up to
$6\cdot10^{-8}$ (single-precision rounding in Ripser).

\paragraph{Cost.} With Ripser, one subsample of $80$ points in $\R^{20}$ takes about
$\CostDescOne$\,s on one CPU core; the pure-NumPy $H_1$ reduction takes about
$\CostHoneNumpy$\,s. The descriptor is computed once per client. Its cost is therefore
negligible compared with training, even with $n_{\mathrm{rep}}=5$ subsamples.
Table~\ref{tab:comm} gives the communication per round.

\begin{table}[h]
\centering\small
\caption{Communication per client and per round (in floats); $p$ is the model dimension.}
\label{tab:comm}
\restable{tab_comm}
\end{table}

% ==========================================================================
\section{Hyper-parameters}
\label{app:hparams}
% ==========================================================================
\begin{table}[h]
\centering\small
\caption{Hyper-parameters, fixed before the experiments and shared by all federations.}
\label{tab:hparams}
\restable{tab_hparams}
\end{table}
\paragraph{Synthetic generators.}
\begin{itemize}[leftmargin=1.2em,itemsep=2pt]
  \item Federation (A): $d=20$; the subspace dimension is $3$; the concept shift has unit
    norm along a direction orthogonal to $\beta_0$; logits are rescaled to standard
    deviation $3$.
  \item Corrupted clients: the shift-plus-flips attack adds $\mathcal N(2,0.5^2)$ to $10$
    coordinates and flips $40\%$ of the labels; the heavy-tailed attack adds Student-$t_2$
    noise to all coordinates and flips $40\%$ of the labels.
  \item Privacy experiment: $n_k\in[400,1200]$, rows rescaled to $\norm{\tilde x}\le C=\DPclip$,
    no pooled standardisation (it would be an additional release), $T=40$, $\tau=5$ and
    $\delta=10^{-5}$.
\end{itemize}

% ==========================================================================
\section{Additional results}
\label{app:extra}
% ==========================================================================

\begin{table}[h]
\centering\small
\caption{\textbf{Per-client view} (all federations). Shown are the fraction of honest
clients whose personalised log-loss is worse than under FedAvg and under local training, and
the loss of the worst client, all averaged over seeds.}
\label{tab:perclient}
\resizebox{\textwidth}{!}{\restable{tab_perclient}}
\end{table}

\begin{table}[h]
\centering\small
\caption{\textbf{Blending and local steps} (federation A, log-loss). Rows: number of local
steps $\tau$, with the total number of local steps held equal. Columns: target effective
blending $\bar\alpha$, with $\alpha$ set by Eq.~\eqref{eq:alphacal}. With calibration, the
results depend on $\bar\alpha$ but hardly on $\tau$, as Proposition~\ref{prop:blend}(c)
predicts.}
\label{tab:alphatau}
\restable{tab_alpha_tau}
\end{table}

\begin{table}[h]
\centering\small
\caption{\textbf{Sensitivity} of \ptopofl\ on federation (A): log-loss and recovery of the
latent groups (ARI).}
\label{tab:sensitivity}
\restable{tab_sensitivity}
\end{table}

\begin{table}[h]
\centering\small
\caption{\textbf{Detection of corrupted clients} by trust weights (federation A, $K=10$).
Each entry gives the fraction of corrupted clients flagged ($t_k<1$) and the fraction of
honest clients flagged, for each fraction of corrupted clients.}
\label{tab:detection}
\restable{tab_detection}
\end{table}

\begin{table}[h]
\centering\small
\caption{\textbf{Privacy experiment} (Section~\ref{sec:e_privacy}): log-loss,
membership-inference AUC, partition recovery and property-inference accuracy as functions
of the record-level $\varepsilon$ ($\delta=10^{-5}$).}
\label{tab:privacy}
\resizebox{\textwidth}{!}{\restable{tab_privacy}}
\end{table}

\paragraph{Drift monitoring.}
When client data change over time, descriptors can be recomputed periodically to monitor
drift. In this experiment, $K=6$ clients receive fresh samples in each of $20$ rounds, and
two of them drift progressively from isotropic to near-$3$-dimensional data. The statistic
$\Delta_k=\frac1R\sum_r\norm{\hat\phi_k^{(r)}-\hat\phi_k^{(1)}}$ separates drifting from
stable clients (AUC $\ContAUC$; flagging the top quartile gives TPR $\ContTPR$ and FPR
$\ContFPR$; Figure~\ref{fig:continual}). With static data, descriptors are computed once,
so no spurious drift arises.

\begin{figure}[h]
\centering
\resfig{fig_continual}{0.8\textwidth}
\caption{Entropy coordinates of the normalised descriptor over $20$ rounds. Dashed lines are
the two drifting clients.}
\label{fig:continual}
\end{figure}

% ==========================================================================
\section{Reproducibility}
\label{app:repro}
% ==========================================================================
All results are produced by one script, \texttt{ptopofl\_v2.py}, which depends only on
NumPy, SciPy, scikit-learn and Matplotlib (Ripser is optional). Running
\texttt{python ptopofl\_v2.py --exp all --seeds 10} executes every experiment with seeds
$0$--$9$. The script writes the \LaTeX{} tables and the numerical macros used in the text,
so that the text and the code cannot diverge. The real dataset is downloaded from the UCI
repository, and the script prints the site sizes after preprocessing. On one CPU core the
synthetic experiments take about $12$ minutes in total.
\codeavailability
